\section{Methodology}\label{sec:method}

The protocol was written to answer the paper's question in full, and our measurement campaign ran part of it. This
section gives both, the design and what actually ran, item by item, so that a reader can tell a measured result from a
designed one and can run the rest with the same tools.

\subsection{Six questions}
\begin{description}
\item[RQ1 (feasibility).] On two 96\,GB GPUs, which configurations of the model serve at all without host offload, and
with which exact failure do the others stop?
\item[RQ2 (cost of offload).] How do reading, writing, TTFT, TPOT and throughput change with the share of experts in
host memory, and how much of that is PCIe traffic?
\item[RQ3 (what the refinement adds).] At equal host bytes, how does Diffbot's per-expert split of hot and cold experts
compare with upstream vLLM's own selective zero-copy offload, and how much depends on the placement policy and its
match to the traffic?
\item[RQ4 (frontier).] What is the cost per token and per concurrent request under pre-registered service-level
objectives, from 1 to 64 users, and how does it compare with a reference that holds everything in GPU memory?
\item[RQ5 (quality).] How much accuracy does the \dv{cfg.exl3_bits}-bit stack lose against the original checkpoint on
reasoning, code, instruction following and long context?
\item[RQ6 (coding agents).] How many simultaneous coding-agent streams keep at least 40 output tokens/s each, and how
many intermittently active developers can share one pod under an observed usage profile (Section~\ref{sec:agent})?
\end{description}

\subsection{An ablation ladder}
To separate what each layer contributes, the protocol climbs a ladder in which each rung adds one layer on the same
hardware and with the same serving settings (\datafile{stack/variants/}, \datafile{experiments/experiments.yaml}).
\begin{enumerate}
\item[(a)] \textbf{upstream}: vLLM at the pinned commit, unpatched, with DeepSeek's original checkpoint. It is expected
to fail at load on two 96\,GB GPUs (Section~\ref{sec:system}), and E2 records the error. \textbf{upstream-uva} is the
same with upstream vLLM's selective zero-copy offload of the expert parameters, the closest upstream baseline that
may fit.
\item[(b)] \textbf{exl3-min} adds the five vLLM patches the pack needs on sm\_120, the unpatched vllm-exl3 plugin and
stock ExLlamaV3 kernels, with every expert in GPU memory. It is expected to fail at load, and E2 records it.
\item[(c0)] \textbf{exl3-uva} is (b) plus upstream vLLM's selective zero-copy offload of the EXL3 expert parameters,
sized to the same host bytes as a plan (\texttt{hostoffload.weights.plan bytes}).
\item[(c)] \textbf{exl3-offload} adds Diffbot's patch, kernels and hooks. Its factors are the plan (host shares from
25 to 60\%, chosen from routing counts of a calibration workload on a disjoint corpus half, the recipe's three plans,
and a seeded random plan with the same sizes per layer), the format of the Engram tables (int4, or FP8 if host memory
allows) and the number of SMs given to the host-expert stream.
\item[(d)] \textbf{DSpark on or off} on top of (c).
\end{enumerate}
A reference rung, \textbf{upstream on 2x B300}, serves the original checkpoint entirely in GPU memory (E9, E11). It is
the baseline of speed and quality that the offload has to be judged against, on different and more expensive
hardware.

\paragraph{The rungs that ran.} Our three pod sessions ran rung (b) and the offload build with an empty plan, both only
to the load failure (E2), rung (c0), which boots or fails but never serves (Section~\ref{sec:results-controlled}), and
rung (c) with DSpark on and two plans: the recipe's \texttt{gen30} plan and the union of its \texttt{gen30} and
\texttt{agent30} cold sets. Rung (a) and its upstream-uva variant (both with the original checkpoint), the calibrated,
random and FP8-Engram plans, the sweep of SM budgets, rung (d) and the B300 reference did not run
(Section~\ref{sec:not-run}).

\subsection{The benchmark protocol}
The harness (\datafile{src/hostoffload/harness/}) replaces the exploratory load generator and fixes the weaknesses
listed in Table~\ref{tab:fixes}. The protocol below is the full design (\datafile{experiments/workloads/main.yaml}),
and each item ends with what our runs used (\datafile{experiments/workloads/lean_main.yaml}, \datafile{experiments/workloads/lean_probes.yaml},
\datafile{experiments/workloads/lean_subset.yaml}, \datafile{experiments/workloads/long_context.yaml} and
\datafile{experiments/workloads/agent.yaml}, run by \datafile{experiments/runs/lean.sh},
\datafile{experiments/runs/long.sh} and \datafile{experiments/runs/agent.sh}).
\begin{itemize}
\item \textbf{Prompts} are cut to exact token counts with the server's own tokenizer endpoints, from pinned corpora:
the vLLM sources at the pinned commit for code, and four public-domain books for prose, both pinned by the sha256 of
the assembled text (\datafile{experiments/corpora.yaml}). Every prompt of a run is unique except deliberately shared
prefixes, and benchmark and calibration workloads use disjoint halves of the corpus. \emph{In our runs} it was as
designed, but every measured prompt came from the code corpus and the calibration workload did not run.
\item \textbf{Scenarios.} A sends independent 4K-token code prompts, A-Poisson the same as an open-loop Poisson stream
at five rates, B a 16K-token shared prefix (cached after one warm request) plus a question, P a 1K-token prose
continuation, where speculation is expected to accept fewer tokens, and L fresh 128K-token prompts. All use 256 output
tokens, temperature 0.6, top-p 0.95, a seed per request and thinking off. \emph{In our runs} only A and B ran in E8,
and neither A-Poisson nor P ran. Long prompts were measured by the separate experiment E12 (fresh prompts up to
260,000 tokens, decode after 128K and 256K-token prompts, and shared prefixes of 128K and 256K tokens), and E13 used
400-token replies (Section~\ref{sec:agent}).
\item \textbf{Arrivals.} Closed loops with $N\in\{1,2,4,8,16,32,64\}$ users whose starts are staggered over 10\,s, so
that there is no synchronized first wave, and Poisson open loops. A sweep stops when TTFT p95 exceeds 120\,s or more
than 5\% of requests fail, and the saturation point is recorded. \emph{In our runs} only closed loops ran. Starts were
staggered over 5\,s in the scenarios with independent short prompts or a shared prefix (E8, E3, the prefix scenarios
of E12 and regimes A and B of E13), but not in the decode scenarios after long prompts of E12 or in regimes C and D
of E13. In E8 and E3 scenario A each user sent two requests and each level had at least eight, and in the decode
probes one per user and at least four. The minimum of two errors in the error rule applies to E12 and E13 only (the
second session's rerun of E3 kept the first session's rule to match its baseline), and the long-context and agent
sweeps stop at a TTFT p95 of 900\,s.
\item \textbf{Reading and writing measured apart.} Prefill probes send unique prompts of 1K to 64K tokens with one
output token, one at a time, so that TTFT is reading plus an empty queue. Decode probes send 128-token prompts with
512 forced output tokens at each $N$, so that TPOT is writing under batching. Server-side histograms of queue, prefill
and decode time per request are differenced around every level. \emph{In our runs} the prefill probes used 4K, 16K and
64K tokens (three requests each), and the decode probes $N\in\{1,8,32\}$, with one repetition.
\item \textbf{Repetitions.} Three per load point within a boot, with fresh prompts and seeds each time, and three cold
boots at $N\in\{1,8,32\}$ for the component that varies between boots. Every request of every repetition is stored
raw. \emph{In our runs} there were no cold-boot repetitions, and each experiment's measurements come from a single
boot. The \texttt{gen30} configuration was booted on three pods, but no measurement was repeated across boots. E8 has
two repetitions per level in B and, after an early stop (Section~\ref{sec:results-controlled}), in A only at $N=1$.
E3, E12 and E13 have one. Every request is stored raw.
\end{itemize}

\begin{table}[t]
\centering\small
\begin{tabular}{p{4.1cm}p{5.6cm}p{4.5cm}}
\toprule
weakness of the exploratory runs & fix in the controlled protocol & in our runs \\
\midrule
single run per configuration & 3 repetitions per point plus 3 cold boots, cluster bootstrap over repetitions, variance components & partly: 1 or 2 repetitions inside one boot, no cold boots \\
usable bars chosen after seeing the data & SLOs registered in \datafile{experiments/sla.yaml} before any controlled run & yes \\
all $N$ requests sent at once (burst) & staggered closed-loop starts, plus open-loop Poisson arrivals & staggered starts (5\,s), no Poisson \\
speculation makes decode content-dependent & separate code and prose scenarios, DSpark on and off (E7), acceptance length recorded & partly: code only, DSpark always on, acceptance length recorded (Table~\ref{tab:spec-acceptance}) \\
quality never evaluated & E10 and E11 with lm-evaluation-harness against the original checkpoint & partly: three task subsets, DeepSeek's API as reference \\
per-chunk timing, several tokens per chunk & cumulative token counts on every chunk, per-token latency weighted by tokens & yes \\
corpus from the installed package, chars-per-token guess & pinned corpora, exact token counts via the server tokenizer & yes \\
5\,s sampler, no PCIe, no per-process memory & 0.5\,s sampler with NVML PCIe RX/TX, cgroup and per-process memory by kind, mapping dump & sampler yes, PCIe per level reported (Table~\ref{tab:pcie-per-token}), no Nsight capture, mappings not analyzed \\
plan of the sweep not recorded & plan path and sha256 in every run's configuration log & yes \\
aggregate rate includes ramp and drain & windowed aggregate rate (first send plus warm-up to last send) & yes \\
\bottomrule
\end{tabular}
\caption{Known weaknesses of the exploratory measurements, how the controlled protocol addresses them, and which fixes
our three pod sessions applied.}
\label{tab:fixes}
\end{table}

\subsection{What we compute}
TTFT is the arrival of the first streamed token minus the send time. TPOT is the time from the first to the last token
divided by the number of tokens after the first, and the per-request rate is its inverse. Per-token latency weights
each chunk's inter-arrival time by the tokens it carried. E2E is the end of the stream minus the send time. Aggregate
output throughput counts the tokens of requests that finished inside a window from the first send (plus a warm-up) to
the last send, so that ramp-up and drain do not dilute it. We report p50, p95 and p99 with linear interpolation, always
with $n$, and with fewer than 20 requests a p95 is close to the maximum. Requests of one repetition share server state,
so uncertainty comes from a bootstrap that resamples whole repetitions~\citep{efron1994bootstrap} and from the
coefficient of variation across repetitions, and a one-way variance decomposition separates the variance between and
within repetitions (\datafile{src/hostoffload/analysis/stats.py}).

\subsection{Following the bytes}
Zero-copy reads never show up as copies, so finding them takes several instruments. For every run we record the memory
of each GPU over time (NVML), the cgroup's anonymous, shared and file memory, each process's resident, anonymous,
shared, locked and pinned memory, the largest mappings of each server process at start, and the placement the server
logs at load (host experts per layer, Engram bytes per rank, weight, KV and graph memory). The static split
(Table~\ref{tab:static-placement}) gives file bytes per tensor class in GPU and host memory, and the runtime records
give allocations. PCIe receive and transmit rates per GPU come from NVML every 0.5\,s and from
\texttt{nvidia-smi dmon}, to be cross-checked by one Nsight Systems capture with GPU metrics and calibrated by a script
that drives a known volume of copy and zero-copy traffic and compares it with the counters
(\datafile{experiments/tools/pcie_calibrate.py}). The analysis correlates the PCIe receive rate with the decode rate
across load points and plans, with Pearson and Spearman coefficients and bootstrap intervals. \emph{In our runs} the
sampler and \texttt{nvidia-smi dmon} ran in every session and the calibration script once, in the first, where its
zero-copy test on the second GPU is invalid. We report PCIe receive traffic per level and per output token
(Section~\ref{sec:bytes}), but the Nsight capture, the attribution of allocations from the mappings and the
correlation analysis were not done.

\subsection{Cost and service-level objectives}
Cost per hour is the pod's quoted price in USD, converted to EUR at the ECB reference rate of \dv{price.fx.date}
(\dv{price.fx.usd_per_eur} USD per EUR, \datafile{experiments/pricing.yaml}). Cost per million output tokens divides
it by the sustained output rate. Splitting cost between input and output tokens is a modelling choice, so the model
offers three rules: all cost to output, a dedicated split (input cost from the prefill probe at saturation, output cost
from the decode probe), and a fixed price ratio solved at a mixed load point. Cost per concurrent request is the hourly
price divided by $N$, and cost per request is the level's pod time divided by its requests. Cost under an objective
uses the highest-throughput load point that meets it in every repetition. The three objectives are interactive (TTFT
p95 up to 2\,s and TPOT p95 up to 50\,ms), relaxed (10\,s and 100\,ms) and long context (5\,s and 66.7\,ms). They
were registered on 2026-10-06 before any controlled run, but after we had seen the exploratory data
(\datafile{experiments/sla.yaml}).

\subsection{Quality}
E10 and E11 run lm-evaluation-harness~\citep{sutawika2023lmeval} 0.4.13 (task names validated locally,
\datafile{quality/tasks.yaml}) against the measured stack and the reference. The design covers GSM8K~\citep{cobbe2021gsm8k}
and MATH-500~\citep{hendrycks2021math} for reasoning, HumanEval~\citep{chen2021humaneval} and MBPP~\citep{austin2021mbpp}
for code, IFEval~\citep{zhou2023ifeval} for instruction following, and subsets of RULER~\citep{hsieh2024ruler} and
LongBench~\citep{bai2024longbench} for long context. Both systems get identical prompts, greedy decoding, thinking off
and fixed seeds, with logs per item, so that differences are tested item by item with a paired bootstrap. Quantization
can change behaviour that averages hide~\citep{dutta2024accuracy}, so we also report the share of items whose outcome
flips. As designed, the reference is DeepSeek's original checkpoint on upstream vLLM on the B300 pod (E11). \emph{In
our runs} only fixed-seed subsets of GSM8K, IFEval and HumanEval were evaluated
(\datafile{quality/subsets_lean2h_seed20261006.json}), and the reference is the official DeepSeek API with thinking
disabled, reached through a local proxy that logs token usage (\datafile{quality/api_proxy.py}), reported separately
because its serving precision and implementation cannot be verified. MATH-500, MBPP, RULER and LongBench did not run.

\subsection{Designed but not run}\label{sec:not-run}
Six experiments of the protocol were designed and implemented but did not run in our campaign. Their full
specifications, workloads and runners are in \datafile{experiments/experiments.yaml} and \datafile{experiments/runs/},
and none of their numbers appears in this paper. E9 would serve the original checkpoint entirely in GPU memory on two
B300 GPUs, the reference against which the cost of the offload would be measured. E7 would turn DSpark off to separate
its effect, on code and on prose. E4 would compare placement policies at a fixed host share, including a random plan
and one fitted to our own routing counts. E5 would compare int4 and FP8 Engram tables in host memory. E6 would vary the
number of SMs given to the stream that computes host experts. E1 would rebuild every variant from scratch on a CPU
machine and record image digests.
